% Part: turing-machines
% Chapter: undecidability
% Section: representing-tms

\documentclass[../../../include/open-logic-section]{subfiles}

\begin{document}

\olfileid{tur}{und}{rep}
\olsection{بازنمایی ماشین‌های تورینگ}

\begin{explain}
برای بازنمایی ماشین‌های تورینگ و رفتارشان با !!a{sentence}ای از منطق
مرتبه‌اول، باید زبان مناسبی تعریف کنیم. این زبان از دو بخش تشکیل
می‌شود: !!{predicate}sیی برای توصیف پیکربندی‌های ماشین، و عبارت‌هایی
برای شماره‌گذاری گام‌های اجرا («لحظه‌ها») و مکان‌های روی نوار.

دو نوع !!{predicate}s معرفی می‌کنیم که هر دو دوموضعی‌اند: به ازای هر
حالت~$q$، !!a{predicate}ی چون~$\Obj Q_q$؛ و به ازای هر نماد نواری
$\sigma$، !!a{predicate}ی چون~$\Obj S_\sigma$. اولی‌ها به ما امکان می‌دهند
حالت~$M$ و مکان سر نوارش را توصیف کنیم و دومی‌ها برای توصیف محتوای
نوار به کار می‌روند.

برای بیان مکان‌های سر نوار و تعداد گام‌های اجراشده به راهی برای بیان
عددها نیاز داریم. این کار با !!a{constant}~$\Obj 0$ و تابع یک‌موضعیِ
$\prime$، یعنی تابع جانشین، انجام می‌شود. بنا بر قرارداد این تابع
\emph{پس از} آرگومانش نوشته می‌شود (و پرانتزها را حذف می‌کنیم).

به ازای هر عدد~$n$، ترم متعارفی چون~$\num{n}$ وجود دارد که
\emph{عددنمای}~$n$ است و آن را در~$\Lang L_M$ بازنمایی می‌کند.
$\num{0}$ همان~$\Obj 0$، $\num{1}$ همان~$\Obj 0'$، $\num{2}$ همان
$\Obj 0''$ و به همین ترتیب است. به‌طور رسمی‌تر:
\begin{align*}
\num{0} & = \Obj 0 \\
\num{n+1} &= \num{n}'
\end{align*}
عبارت~$\num{0}$، یعنی~$\Obj 0$، هم چپ‌ترین مکان روی نوار و هم زمان
پیش از نخستین گام اجرا (پیکربندی آغازین) را نام‌گذاری می‌کند. عبارت
$\num{1}$، یعنی~$\Obj 0'$، خانهٔ سمت راست چپ‌ترین خانه و زمان پس از
نخستین گام اجرا را نام‌گذاری می‌کند و به همین ترتیب.

همچنین !!a{predicate}~$<$ را معرفی می‌کنیم تا هم ترتیب مکان‌های نوار (که
در آن صورت به معنای «در سمت چپِ» است) و هم ترتیب گام‌های اجرا (که در
آن صورت به معنای «پیش از» است) را بیان کنیم.

پس از فراهم‌کردن زبان، «اصول موضوعه»ٔ~$!T(M,w)$ را فهرست می‌کنیم؛
یعنی !!{sentence}sیی که در کنار هم رفتار~$M$ را هنگام اجرا روی ورودی~$w$
توصیف می‌کنند. !!{sentence}sیی شرایط مربوط به~$\Obj 0$، $\prime$ و~$<$ را
تعیین می‌کنند؛ !!{sentence}sیی پیکربندی ورودی را توصیف می‌کنند؛ و !!{sentence}sیی
پیکربندی~$M$ را پس از اجرای یک دستور خاص شرح می‌دهند.
\end{explain}

\begin{defn}
  \ollabel{defn:tm-descr}
اگر ماشین تورینگ $M=\tuple{Q,\Sigma,q_0,\delta}$ داده شده باشد، زبان
$\Lang L_M$ شامل موارد زیر است:
\begin{enumerate}
\item به ازای هر حالت~$q\in Q$ یک !!{predicate} دوموضعیِ
  $\Obj Q_q(x,y)$. به‌طور شهودی، $\Obj Q_q(\num{m},\num{n})$ بیان
  می‌کند «پس از $n$ گام، $M$ در حالت~$q$ است و خانهٔ $m$اُم را پویش
  می‌کند».
\item به ازای هر نماد~$\sigma\in\Sigma$ یک !!{predicate} دوموضعیِ
  $\Obj S_\sigma(x,y)$. به‌طور شهودی،
  $\Obj S_\sigma(\num{m},\num{n})$ بیان می‌کند «پس از $n$ گام، خانهٔ
  $m$اُم شامل نماد~$\sigma$ است».
\item یک !!{constant}~$\Obj 0$؛
\item یک !!{function} یک‌موضعیِ~$\prime$؛
\item یک !!{predicate} دوموضعیِ~$<$.
\end{enumerate}
\end{defn}

!!{sentence}sیی که کار ماشین تورینگ~$M$ را روی ورودی
$w=\sigma_{i_1}\dots\sigma_{i_k}$ توصیف می‌کنند از این قرارند:
\begin{enumerate}
\item اصول موضوعه‌ای که اعداد و~$<$ را توصیف می‌کنند:
\begin{enumerate}
\item !!^a{sentence} که می‌گوید هر عدد از جانشینش کوچک‌تر است:
\[
\lforall[x][x < x']
\]
\item !!^a{sentence} که متعدی‌بودن~$<$ را تضمین می‌کند:
\[
\lforall[x][\lforall[y][\lforall[z][
      ((x < y \land y < z) \lif x < z)]]]
\]
\end{enumerate}
\item اصول موضوعه‌ای که پیکربندی ورودی را توصیف می‌کنند:
\begin{enumerate}
\item پس از $0$ گام، یعنی پیش از آغاز ماشین، $M$ در حالت آغازین
  $q_0$ است و خانهٔ~$1$ را پویش می‌کند:
\[
\Obj Q_{q_0}(\num{1}, \num{0})
\]
\item نخستین $k+1$ خانه شامل نمادهای~$\TMendtape$،
  $\sigma_{i_1}$،~\dots، $\sigma_{i_k}$ است:
\[
\Obj S_\TMendtape(\num{0}, \num{0}) \land
\Obj S_{\sigma_{i_1}}(\num{1}, \num{0}) \land
\dots \land
\Obj S_{\sigma_{i_k}}(\num{k}, \num{0})
\]
\item نوار در جاهای دیگر خالی است:
\[
\lforall[x][(\num{k} < x \lif \Obj S_\TMblank(x, \num{0}))]
\]
\end{enumerate}
\item اصول موضوعه‌ای که گذار از یک پیکربندی به پیکربندی بعدی را
  توصیف می‌کنند:

در ادامه، $!A(x,y)$ را عطف همهٔ !!{sentence}sیی از صورت زیر بگیرید:
\[
\lforall[z][
  (((z < x \lor x < z) \land \Obj S_\sigma(z, y))
  \lif \Obj S_\sigma(z, y'))]
\]
که در آن~$\sigma\in\Sigma$. از~$!A(\num{m},\num{n})$ استفاده می‌کنیم
تا بیان کنیم «جز در خانهٔ~$m$، نوار پس از $n+1$ گام همان است که پس
از $n$ گام بود».
\begin{enumerate}
\item \ollabel{rep-right} به ازای هر دستور
  $\delta(q_i,\sigma)=\tuple{q_j,\sigma',\TMright}$، !!{sentence}ٔ زیر:
\begin{align*}
& \lforall[x][\lforall[y][(
   (\Obj Q_{q_i}(x, y) \land \Obj S_{\sigma}(x, y)) \lif {}]] \\
&\qquad   (\Obj Q_{q_j}(x', y') \land \Obj S_{\sigma'}(x, y') \land
!A(x, y)))
\end{align*}
این می‌گوید اگر پس از~$y$ گام ماشین در حالت~$q_i$ باشد و خانهٔ~$x$
دارای نماد~$\sigma$ را پویش کند، پس از $y+1$ گام خانهٔ~$x+1$ را پویش
می‌کند، در حالت~$q_j$ است، خانهٔ~$x$ اکنون شامل~$\sigma'$ است، و هر
خانه‌ای جز~$x$ همان نمادی را دارد که پس از~$y$ گام داشت.

\item \ollabel{rep-left} به ازای هر دستور
  $\delta(q_i,\sigma)=\tuple{q_j,\sigma',\TMleft}$، !!{sentence}ٔ زیر:
\begin{align*}
& \lforall[x][\lforall[y][
    ((\Obj Q_{q_i}(x', y) \land \Obj S_{\sigma}(x', y)) \lif {}]]\\
& \qquad   (\Obj Q_{q_j}(x, y') \land \Obj S_{\sigma'}(x', y') \land
!A(x', y))) \land {}\\
& \lforall[y][((\Obj Q_{q_i}(\num{0}, y) \land \Obj S_{\sigma}(\num{0},
    y)) \lif {}]\\
& \qquad (\Obj Q_{q_j}(\num{0}, y') \land \Obj S_{\sigma'}(\num{0},
  y') \land !A(\num{0}, y)))
\end{align*}
لحظه‌ای دربارهٔ چگونگی کار این فرمول بیندیشید: این بار با «اگر خانهٔ
$x$ را پویش کند~\dots» آغاز نمی‌کنیم، بلکه می‌گوییم «اگر خانهٔ $x+1$
را پویش کند~\dots». حرکت به چپ یعنی ماشین در گام بعد خانهٔ~$x$ را
پویش می‌کند. اما خانه‌ای که روی آن نوشته می‌شود~$x+1$ است. این کار را
چنین انجام می‌دهیم زیرا تفریق یا تابع پیشین نداریم.

توجه کنید اعداد به صورت~$x+1$ عبارت‌اند از~$1$، $2$،~\dots؛ پس این
حالت را پوشش نمی‌دهد که ماشین خانهٔ~$0$ را پویش می‌کند و باید به چپ
حرکت کند (که البته نمی‌تواند و همان‌جا می‌ماند). عطف دوم این حالت
ویژه را پوشش می‌دهد: می‌گوید اگر پس از~$y$ گام ماشین در حالت~$q_i$
خانهٔ~$0$ را پویش کند و خانهٔ~$0$ شامل نماد~$\sigma$ باشد، پس از
$y+1$ گام همچنان خانهٔ~$0$ را پویش می‌کند، اکنون در حالت~$q_j$ است،
نماد خانهٔ~$0$ برابر~$\sigma'$ است، و خانه‌های دیگر همان نمادهایی را
دارند که پس از~$y$ گام داشتند.
\item \ollabel{rep-stay} به ازای هر دستور
  $\delta(q_i,\sigma)=\tuple{q_j,\sigma',\TMstay}$، !!{sentence}ٔ زیر:
\begin{align*}
& \lforall[x][\lforall[y][(
   (\Obj Q_{q_i}(x, y) \land \Obj S_{\sigma}(x, y)) \lif {}]] \\
&\qquad   (\Obj Q_{q_j}(x, y') \land \Obj S_{\sigma'}(x, y') \land
!A(x, y)))
\end{align*}
\end{enumerate}
\end{enumerate}
بگذارید $!T(M,w)$ عطف همهٔ !!{sentence}sی بالا برای ماشین تورینگ~$M$ و
ورودی~$w$ باشد.

برای بیان اینکه $M$ سرانجام متوقف می‌شود، باید !!{sentence}ای بیابیم که
بگوید «پس از تعدادی گام، تابع گذار تعریف‌نشده خواهد بود». بگذارید~$X$
مجموعهٔ همهٔ زوج‌های~$\tuple{q,\sigma}$ باشد که
$\delta(q,\sigma)$ برایشان تعریف‌نشده است. آنگاه بگذارید $!E(M,w)$
!!{sentence}ٔ زیر باشد:
\[
\lexists[x][\lexists[y][(\bigvee_{\tuple{q, \sigma} \in
      X}(\Obj Q_q(x, y) \land \Obj S_\sigma(x, y)))]]
\]

اگر ماشین تورینگی با حالت توقف معین~$h$ به کار ببریم، حتی ساده‌تر
است: در آن صورت !!{sentence}ٔ~$!E(M,w)$
\[
\lexists[x][\lexists[y][\Obj Q_h(x, y)]]
\]
بیان می‌کند که ماشین سرانجام متوقف می‌شود.

\begin{prop}
\ollabel{prop:mlessk}
اگر $m<k$، آنگاه $!T(M,w)\Entails\num{m}<\num{k}$.
\end{prop}

\begin{proof}
تمرین.
\end{proof}

\begin{prob}
گزارهٔ \olref[tur][und][rep]{prop:mlessk} را ثابت کنید.
(راهنمایی: بر~$k-m$ استقرا کنید.)
\end{prob}

\end{document}
